\ifajar
  \section{Latihan}
\else
  \section{Bab 8: Arsitektur Web Modern}
\fi

Seluruh file yang dipakai pada latihan ini tersedia di
\berkas{sources/chapter08/}, dihitung relatif terhadap akar repositori.
Latihan ini hanya membutuhkan Python 3, tanpa Docker dan tanpa Node.js,
karena yang dibandingkan adalah tempat HTML di-\textit{render}.

Seluruh pengukuran diarahkan ke aplikasi yang berjalan di komputer sendiri,
dengan beban setara satu pengunjung membuka halaman. Hasil tiap mahasiswa
akan berbeda pada bagian waktu, dan yang dilaporkan adalah pola
perbandingannya beserta penjelasannya.

Penyiapannya dilakukan sekali, seperti pada Listing \ref{lst:c8-siapkan}.

\begin{lstlisting}[language=bash, caption={Membangun halaman statis lalu menyalakan server, memakai \berkas{sources/chapter08/bangun-statis.py}}, label={lst:c8-siapkan}]
cd sources/chapter08
python3 bangun-statis.py
python3 aplikasi.py &

# Keluarannya:
# Dibangun     : katalog-ssg.html
# Ukuran       : 511 byte
# Lama render  : 0.5 ms
# Melayani di http://127.0.0.1:8040/ssr  (Ctrl+C untuk berhenti)
\end{lstlisting}

\ifajar
  \subsection*{Latihan 1: Membandingkan Lima Varian}
\else
  \subsection*{Prosedur 1: Membandingkan Lima Varian}
\fi

Latihan ini membuktikan bahwa perbedaan kelima varian dapat dihitung, bukan
hanya diperdebatkan.

\begin{enumerate}
  \item Jalankan alat ukurnya seperti pada Listing \ref{lst:c8-lat1-jalan},
        lalu isi Tabel \ref{tab:c8-lat1-hasil}. Tabel
        \ref{tab:c8-lat1-contoh} memperlihatkan contoh yang sudah terisi.
  \item Buka kelima alamat di browser, lalu bandingkan apa yang tampil
        ketika JavaScript dimatikan lewat panel pengembang.
  \item Lihat sumber halaman tiap varian, lalu catat varian mana yang
        memuat judul buku di dalam HTML-nya.
  \item Jawab tiga pertanyaan berikut. Pertama, mengapa HTML varian SPA
        paling kecil tetapi totalnya paling besar? Kedua, mengapa alat ukur
        ini tidak dapat menangkap seluruh biaya varian SPA? Ketiga, varian
        mana yang paling tahan terhadap skrip yang gagal dimuat?
\end{enumerate}

\begin{lstlisting}[language=bash, caption={Mengukur kelima varian, tersedia di \berkas{sources/chapter08/ukur-varian.py}}, label={lst:c8-lat1-jalan}]
python3 ukur-varian.py

# Keluarannya:
#   Varian        HTML   Total  Isi di HTML  Minta  HTML ms  Total ms
#   ssg            511     511           ya      1      1.9       1.9
#   isr            511     511           ya      1     52.9      52.9
#   ssr            511     511           ya      1     51.8      51.8
#   islands        625    1322           ya      2      1.5       3.1
#   spa            229    1551        tidak      3      0.9       1.7
\end{lstlisting}

\begin{table}[!htb]
  \centering
  \caption{Contoh hasil Latihan 1 yang sudah terisi}
  \label{tab:c8-lat1-contoh}
  \begin{tabularx}{\textwidth}{lrrlr}
    \hline
    \textbf{Varian} & \textbf{HTML} & \textbf{Total} & \textbf{Isi di HTML} & \textbf{Minta} \\
    \hline
    SSG & 511 B & 511 B & ya & 1 \\
    \hline
    SPA & 229 B & 1.551 B & tidak & 3 \\
    \hline
  \end{tabularx}
\end{table}

\begin{table}[!htb]
  \centering
  \caption{Lembar isian hasil pengukuran Latihan 1}
  \label{tab:c8-lat1-hasil}
  \begin{tabularx}{\textwidth}{lrrlr}
    \hline
    \textbf{Varian} & \textbf{HTML} & \textbf{Total} & \textbf{Isi di HTML} & \textbf{Minta} \\
    \hline
    SSG & \dots & \dots & \dots & \dots \\
    \hline
    ISR & \dots & \dots & \dots & \dots \\
    \hline
    SSR & \dots & \dots & \dots & \dots \\
    \hline
    Islands & \dots & \dots & \dots & \dots \\
    \hline
    SPA & \dots & \dots & \dots & \dots \\
    \hline
  \end{tabularx}
\end{table}

\ifajar
  \subsection*{Latihan 2: Membaca Perilaku ISR}
\else
  \subsection*{Prosedur 2: Membaca Perilaku ISR}
\fi

Latihan ini membuktikan bahwa biaya \textit{rendering} pada ISR tidak hilang,
melainkan ditanggung satu pengunjung tiap periode.

\begin{enumerate}
  \item Minta alamat ISR tiga kali berurutan, lalu catat header
        \texttt{X-ISR} dan waktunya, seperti pada Listing
        \ref{lst:c8-lat2-jalan}.
  \item Tunggu sebelas detik, lalu minta dua kali lagi dan catat kembali.
  \item Ubah \texttt{UMUR\_ISR\_DETIK} di
        \berkas{sources/chapter08/aplikasi.py} menjadi 60, nyalakan ulang
        aplikasinya, lalu ulangi pengamatannya.
  \item Isi Tabel \ref{tab:c8-lat2-hasil}. Tabel \ref{tab:c8-lat2-contoh}
        memperlihatkan contoh yang sudah terisi.
  \item Jawab dua pertanyaan berikut. Pertama, pengunjung mana yang
        menanggung biaya \textit{rendering}, dan apakah itu adil? Kedua, apa yang
        berubah bila umurnya dinaikkan menjadi satu jam?
\end{enumerate}

\begin{lstlisting}[language=bash, caption={Mengamati daur ISR, tersedia di \berkas{sources/chapter08/aplikasi.py}}, label={lst:c8-lat2-jalan}]
curl -s -D - -o /dev/null -w "%{time_total}\n" http://127.0.0.1:8040/isr \
  | grep -iE "^(x-isr|[0-9])"

# Keluarannya:
# X-ISR: di-render ulang
# 0.052217
\end{lstlisting}

\begin{table}[!htb]
  \centering
  \caption{Contoh hasil Latihan 2 yang sudah terisi, umur file sepuluh detik}
  \label{tab:c8-lat2-contoh}
  \begin{tabularx}{\textwidth}{rXr}
    \hline
    \textbf{Permintaan} & \textbf{Keterangan \texttt{X-ISR}} & \textbf{Waktu} \\
    \hline
    1 & \texttt{di-render ulang} & 52,2 ms \\
    \hline
    2 & file lama dipakai & 1,2 ms \\
    \hline
    4 & \texttt{di-render ulang}, sesudah jeda & 52,6 ms \\
    \hline
  \end{tabularx}
\end{table}

\begin{table}[!htb]
  \centering
  \caption{Lembar isian hasil pengamatan Latihan 2}
  \label{tab:c8-lat2-hasil}
  \begin{tabularx}{\textwidth}{rXr}
    \hline
    \textbf{Permintaan} & \textbf{Keterangan \texttt{X-ISR}} & \textbf{Waktu} \\
    \hline
    1 & \dots & \dots \\
    \hline
    2 & \dots & \dots \\
    \hline
    3 & \dots & \dots \\
    \hline
    4, sesudah jeda sebelas detik & \dots & \dots \\
    \hline
    5 & \dots & \dots \\
    \hline
  \end{tabularx}
\end{table}

\ifajar
  \subsection*{Latihan 3: Memindahkan Satu Bagian Menjadi Pulau}
\else
  \subsection*{Prosedur 3: Memindahkan Satu Bagian Menjadi Pulau}
\fi

Latihan ini membuktikan bahwa jumlah JavaScript yang diunduh ditentukan oleh
seberapa banyak bagian halaman yang dihidupkan.

\begin{enumerate}
  \item Baca \berkas{sources/chapter08/statis/pulau.js}, lalu catat
        ukurannya dan apa saja yang dikerjakannya.
  \item Bandingkan dengan \berkas{sources/chapter08/statis/spa.js}, lalu
        catat ukurannya dan apa yang akan hilang bila skrip itu gagal
        dimuat.
  \item Tambahkan satu pulau kedua pada halaman islands, misalnya tombol
        yang mengurutkan daftar menurut harga, lalu ukur ulang kedua varian.
  \item Isi Tabel \ref{tab:c8-lat3-hasil}. Tabel \ref{tab:c8-lat3-contoh}
        memperlihatkan contoh yang sudah terisi.
  \item Jawab dua pertanyaan berikut. Pertama, berapa byte JavaScript yang
        bertambah untuk tiap pulau baru? Kedua, pada titik mana menambah
        pulau menjadi tidak lagi lebih murah daripada memakai SPA?
\end{enumerate}

\begin{lstlisting}[language=bash, caption={Membandingkan ukuran skrip kedua varian, memakai \berkas{sources/chapter08/statis/}}, label={lst:c8-lat3-jalan}]
wc -c statis/pulau.js statis/spa.js

# Keluarannya:
#  697 statis/pulau.js
#  938 statis/spa.js
# 1635 total
\end{lstlisting}

\begin{table}[!htb]
  \centering
  \caption{Contoh hasil Latihan 3 yang sudah terisi}
  \label{tab:c8-lat3-contoh}
  \begin{tabularx}{\textwidth}{Xrr}
    \hline
    \textbf{Varian} & \textbf{Byte skrip} & \textbf{Byte total} \\
    \hline
    Islands, satu pulau & 697 & 1.322 \\
    \hline
    SPA & 938 & 1.551 \\
    \hline
  \end{tabularx}
\end{table}

\begin{table}[!htb]
  \centering
  \caption{Lembar isian hasil pengukuran Latihan 3}
  \label{tab:c8-lat3-hasil}
  \begin{tabularx}{\textwidth}{Xrr}
    \hline
    \textbf{Varian} & \textbf{Byte skrip} & \textbf{Byte total} \\
    \hline
    Islands, satu pulau & \dots & \dots \\
    \hline
    Islands, dua pulau & \dots & \dots \\
    \hline
    SPA & \dots & \dots \\
    \hline
  \end{tabularx}
\end{table}

\ifajar
  \subsection*{Latihan 4: Menulis Catatan Keputusan}
\else
  \subsection*{Prosedur 4: Menulis Catatan Keputusan}
\fi

Latihan ini mengubah angka pada ketiga latihan sebelumnya menjadi satu
keputusan yang dapat dipertanggungjawabkan.

\begin{enumerate}
  \item Pilih satu kebutuhan nyata, misalnya halaman katalog publik, halaman
        rincian pesanan milik pengguna yang sudah masuk, atau \textit{dashboard}
        admin.
  \item Pilih tiga varian yang masuk akal bagi kebutuhan itu, lalu tulis
        angka pembandingnya dari Tabel \ref{tab:c8-lat1-hasil}.
  \item Tulis catatan keputusan berisi lima bagian seperti pada Tabel
        \ref{tab:c8-lat4-hasil}, yaitu keputusan, pilihan yang
        dipertimbangkan, alasan, konsekuensi, dan syarat yang
        membatalkannya. Tabel \ref{tab:c8-lat4-contoh} memperlihatkan
        contoh yang sudah terisi.
  \item Jawab dua pertanyaan berikut. Pertama, angka mana yang paling
        menentukan pilihan itu, dan mengapa? Kedua, keadaan apa yang akan
        membuat keputusan itu perlu ditinjau ulang?
\end{enumerate}

\begin{table}[!htb]
  \centering
  \caption{Contoh catatan keputusan yang sudah terisi}
  \label{tab:c8-lat4-contoh}
  \begin{tabularx}{\textwidth}{lX}
    \hline
    \textbf{Bagian} & \textbf{Isinya} \\
    \hline
    Keputusan & Halaman katalog publik memakai ISR dengan umur sepuluh menit \\
    \hline
    Dipertimbangkan & SSR, SSG, ISR \\
    \hline
    Alasan & SSR menanggung 51,8 ms tiap permintaan, SSG terlalu kaku untuk isi yang berubah beberapa kali sehari \\
    \hline
    Konsekuensi & Satu pengunjung tiap periode menanggung 52,9 ms, isi dapat \textit{stale} sampai sepuluh menit \\
    \hline
    Yang membatalkan & Isi wajib \textit{fresh} seketika, atau \textit{rendering}-nya menjadi jauh lebih mahal \\
    \hline
  \end{tabularx}
\end{table}

\begin{table}[!htb]
  \centering
  \caption{Lembar isian catatan keputusan Latihan 4}
  \label{tab:c8-lat4-hasil}
  \begin{tabularx}{\textwidth}{lX}
    \hline
    \textbf{Bagian} & \textbf{Isinya} \\
    \hline
    Kebutuhan yang dipilih & \dots \\
    \hline
    Keputusan & \dots \\
    \hline
    Dipertimbangkan & \dots \\
    \hline
    Alasan, beserta angkanya & \dots \\
    \hline
    Konsekuensi & \dots \\
    \hline
    Yang membatalkan & \dots \\
    \hline
  \end{tabularx}
\end{table}

Setelah seluruh latihan selesai, server dihentikan dengan Ctrl+C. Tidak ada
container yang perlu dimatikan, karena bab ini memang tidak memakainya.
